\problem{4.1}{Canonical commutators and three-dimensional Ehrenfest theorem}
For smooth test functions $f(\mathbf r)$,
\begin{align}
[r_i,p_j]f&=-\ii\sh r_i\partial_j f+
\ii\sh\partial_j(r_if)=\ii\sh\delta_{ij}f,\\
[r_i,r_j]&=0,&[p_i,p_j]&=0.
\end{align}
Thus $[r_i,p_j]=\ii\sh\delta_{ij}$, including
$[x,p_y]=0$ and $[x,p_x]=\ii\sh$.
For $H=\sum_jp_j^2/(2m)+V$,
$[H,r_i]=-\ii\sh p_i/m$ and
$[H,p_i]=\ii\sh\partial_iV$. Apply
$\dd\langle Q\rangle/\dd t=(\ii/\sh)
\langle[H,Q]\rangle$:
\begin{equation}
\boxed{\frac{\dd\langle\mathbf r\rangle}{\dd t}
=\frac{\langle\mathbf p\rangle}{m},
\qquad
\frac{\dd\langle\mathbf p\rangle}{\dd t}
=-\langle\nabla V\rangle}.
\end{equation}

\problem{4.2}{Infinite cubic well and accidental degeneracy}
Inside $0<x,y,z<a$, separate
$\psi=X(x)Y(y)Z(z)$, each factor satisfying
a one-dimensional infinite-well equation.
The normalized stationary states and energies are
\begin{align}
\boxed{\psi_{n_xn_yn_z}(\mathbf r)}
&=\boxed{\left(\frac2a\right)^{3/2}
\sin\frac{n_x\pi x}{a}
\sin\frac{n_y\pi y}{a}
\sin\frac{n_z\pi z}{a}},\\
\boxed{E_{n_xn_yn_z}}&=
\frac{\pi^2\sh^2}{2ma^2}(n_x^2+n_y^2+n_z^2).
\end{align}
Writing $E_*=\pi^2\sh^2/(2ma^2)$, distinct
levels begin
\begin{equation}
\begin{array}{c|rrrrrr}
\text{level}&E_1&E_2&E_3&E_4&E_5&E_6\\\hline
E/E_*&3&6&9&11&12&14\\
\text{degeneracy}&1&3&3&3&1&6
\end{array}.
\end{equation}
The fourteenth distinct level is
$\boxed{E_{14}=27E_*}$ with
$\boxed{4}$ states: three permutations of $(1,1,5)$
and the distinct state $(3,3,3)$.
The latter is an accidental degeneracy,
not explained by permuting the three axes.

\problem{4.3}{Reconstructing central potentials from wavefunctions}
For a real radial wavefunction $\psi(r)$,
$\nabla^2\psi=\psi''+(2/r)\psi'$, so the
stationary equation gives
\begin{equation}
\boxed{V(r)=E+\frac{\sh^2}{2m}
\frac{\psi''+2\psi'/r}{\psi}}.
\end{equation}
(a) If $\psi=A\ee^{-r/a}$,
$\psi'/\psi=-1/a$,
$\psi''/\psi=1/a^2$, and
\begin{equation}
V(r)=E+\frac{\sh^2}{2m}
\left(\frac1{a^2}-\frac2{ar}\right).
\end{equation}
As $V(\infty)=0$,
$\boxed{E=-\sh^2/(2ma^2)}$ and
$\boxed{V=-\sh^2/(mar)}$.
(b) If $\psi=A\ee^{-r^2/a^2}$,
$\psi'/\psi=-2r/a^2$,
$\psi''/\psi=-2/a^2+4r^2/a^4$, yielding
$\nabla^2\psi/\psi=-6/a^2+4r^2/a^4$.
Setting $V(0)=0$ gives
\begin{equation}
\boxed{E=3\sh^2/(ma^2),\qquad
V(r)=2\sh^2r^2/(ma^4)}.
\end{equation}
This is an isotropic oscillator of frequency
$\omega=2\sh/(ma^2)$.

\problem{4.4}{Constructing and normalizing spherical harmonics}
With the Condon--Shortley convention,
$P_0=1$, $P_2^1(x)=-3x\sqrt{1-x^2}$ and
\begin{equation}
\boxed{Y_0^0=\frac1{\sqrt{4\pi}},
\qquad Y_2^1=-\sqrt{\frac{15}{8\pi}}\sin\theta
\cos\theta\,\ee^{\ii\phi}}.
\end{equation}
For $Y_{00}$,
$\int|Y_{00}|^2\dd\Omega=4\pi/(4\pi)=1$.
For $Y_{21}$,
$\int_0^{2\pi}\dd\phi=2\pi$ and
$\int_0^\pi\sin^3\theta\cos^2\theta\dd\theta
=4/15$, hence its norm is
$(15/(8\pi))(2\pi)(4/15)=1$.
Orthogonality follows by integrating the
factor $\ee^{\ii\phi}$ over $\phi$.

\problem{4.5}{Why the second angular solution is unacceptable}
For $\ell=m=0$ the polar angular equation is
$\frac1{\sin\theta}\frac{\dd}{\dd\theta}
(\sin\theta\,\Theta')=0$. Integrating gives
$\sin\theta\,\Theta'=A$, so
\begin{equation}
\Theta(\theta)=A\ln\tan\frac\theta2+B.
\end{equation}
Although the logarithm is integrable with
the $\sin\theta$ measure, it diverges at
$\theta=0,\pi$, and is not a regular smooth
angular eigenfunction on the sphere.
It is excluded by the regularity domain of
the self-adjoint angular Laplacian; the
acceptable $\ell=0$ solution is constant.
\note{Divergence at a coordinate pole,
not failure of square integrability alone,
is the key domain condition.}

\problem{4.6}{Conjugation identity of spherical harmonics}
Using the real associated Legendre function
and standard normalization,
$Y_\ell^m=N_{\ell m}P_\ell^m(\cos\theta)\ee^{\ii m\phi}$.
For integer $m>0$,
\begin{equation}
P_\ell^{-m}(x)=(-1)^m
\frac{(\ell-m)!}{(\ell+m)!}P_\ell^m(x),\quad
N_{\ell,-m}=
\sqrt{\frac{2\ell+1}{4\pi}
\frac{(\ell+m)!}{(\ell-m)!}}.
\end{equation}
Multiplying these and reversing the
azimuthal exponential gives
\begin{equation}
\boxed{Y_\ell^{-m}=(-1)^m(Y_\ell^{m})^*}.
\end{equation}

\problem{4.7}{Highest-weight harmonics and the explicit $Y_3^2$}
Rodrigues' expression gives
$P_\ell^\ell(x)=(-1)^\ell(2\ell-1)!!
(1-x^2)^{\ell/2}$.
Therefore
\begin{equation}
\boxed{Y_\ell^\ell=
(-1)^\ell
\sqrt{\frac{2\ell+1}{4\pi}
\frac{(2\ell-1)!!^2}{(2\ell)!}}\,
\sin^\ell\theta\,\ee^{\ii\ell\phi}}.
\end{equation}
For $\ell=3$, $P_3^2(x)=15x(1-x^2)$,
and
\begin{equation}
\boxed{Y_3^2=
\sqrt{\frac{105}{32\pi}}\,
\sin^2\theta\cos\theta\,\ee^{2\ii\phi}}.
\end{equation}
Substitution into the angular Laplacian
$\Delta_\Omega=
(\sin\theta)^{-1}\partial_\theta
(\sin\theta\,\partial_\theta)
+(\sin^2\theta)^{-1}\partial_\phi^2$
gives $\Delta_\Omega Y_\ell^m
=-\ell(\ell+1)Y_\ell^m$;
in the second case the eigenvalue is $-12$.

\problem{4.8}{Legendre orthogonality from Rodrigues' formula}
The Legendre functions obey
$\frac{\dd}{\dd x}[(1-x^2)P_\ell']
+\ell(\ell+1)P_\ell=0$.
Multiply the $\ell$ equation by $P_j$,
the $j$ equation by $P_\ell$, subtract,
and integrate over $[-1,1]$:
\begin{equation}
[\ell(\ell+1)-j(j+1)]
\int_{-1}^1P_\ell P_j\dd x
=\left[(1-x^2)(P_jP_\ell'-P_\ell P_j')\right]_{-1}^1=0.
\end{equation}
For $\ell\ne j$ the overlap vanishes.
To normalize, Rodrigues' formula gives
$P_\ell=(2^\ell\ell!)^{-1}\frac{\dd^\ell}{\dd x^\ell}
(x^2-1)^\ell$. Integrating by parts $\ell$
times, boundary terms vanish, leaving
\begin{equation}
\int_{-1}^1P_\ell^2\dd x
=\frac{(2\ell)!}{2^{2\ell}(\ell!)^2}
\int_{-1}^1(1-x^2)^\ell\dd x
=\frac{2}{2\ell+1}.
\end{equation}
Hence
$\boxed{\int_{-1}^1P_\ell P_j\dd x=
2\delta_{\ell j}/(2\ell+1)}$.

\problem{4.9}{Spherical Neumann functions near the origin}
The irregular spherical Bessel functions
are $n_0(x)=-\cos x/x$, $n_1(x)=-\cos x/x^2-\sin x/x$
and $n_{\ell+1}=(2\ell+1)n_\ell/x-n_{\ell-1}$.
Therefore
\begin{equation}
\boxed{n_1(x)=-\frac{\cos x}{x^2}
-\frac{\sin x}{x},\quad
n_2(x)=
-\left(\frac3{x^3}-\frac1x\right)\cos x
-\frac{3\sin x}{x^2}}.
\end{equation}
With $\sin x=x-x^3/6+\cdots$ and
$\cos x=1-x^2/2+\cdots$,
\begin{equation}
\boxed{n_1(x)=-x^{-2}-\tfrac12+O(x^2),\qquad
n_2(x)=-3x^{-3}-\tfrac1{2x}+O(x)}.
\end{equation}
The irregular solutions diverge at $r=0$;
for a finite central potential and
physical regularity one selects $j_\ell$.

\problem{4.10}{Spherical hard-wall energies at $\ell=1$}
For free radial motion $R(r)=Aj_1(kr)$,
$j_1(z)=\sin z/z^2-\cos z/z$.
It obeys $R''+(2/r)R'+[k^2-2/r^2]R=0$
by substitution (the $\ell=1$ radial equation).
A hard wall at $r=a$ imposes $j_1(ka)=0$:
\begin{equation}
\sin z-z\cos z=0
\quad\Longleftrightarrow\quad
\boxed{\tan z=z},\qquad z=ka>0.
\end{equation}
The first three positive roots are
$z_1\simeq4.493409$, $z_2\simeq7.725252$,
$z_3\simeq10.904122$, giving
$\boxed{E_{N1}=\sh^2 z_N^2/(2ma^2)}$.
For large $N$, $z_N\sim(N+1/2)\pi$ and
\begin{equation}
\boxed{E_{N1}\sim\frac{\pi^2\sh^2}{2ma^2}
(N+\tfrac12)^2}.
\end{equation}
The $\ell=1$ magnetic degeneracy is three.

\problem{4.11}{Finite spherical well and the threshold for binding}
For $V=-V_0$ inside $r<a$ and $V=0$
outside, write $u=rR$, $E=-\sh^2\kappa^2/(2m)$,
$k^2=2m(V_0+E)/\sh^2$. For $\ell=0$,
regularity at the origin gives
$u_{\rm in}=A\sin kr$, while
$u_{\rm out}=B\ee^{-\kappa r}$.
Matching $u,u'$ at $a$ gives
\begin{equation}
\boxed{k\cot(ka)=-\kappa},\qquad
k^2+\kappa^2=\frac{2mV_0}{\sh^2}.
\end{equation}
At the first threshold $\kappa\to0^+$,
$ka\to\pi/2$, so a bound level exists
only when $ka$ can reach at least $\pi/2$:
\begin{equation}
\boxed{V_0a^2>\frac{\pi^2\sh^2}{8m}}.
\end{equation}
At equality the zero-energy threshold
state is not square-integrable outside;
there is no normalizable bound state.
This contrasts with one-dimensional
finite wells, which bind at arbitrarily
small integrated attraction.
